\documentclass[11pt,a4paper]{article}
\usepackage[utf8]{inputenc}
\usepackage[T1]{fontenc}
\usepackage{amsmath,amssymb,amsthm}
\usepackage{graphicx}
\usepackage{hyperref}
\usepackage{xcolor}
\usepackage{listings}
\usepackage{booktabs}
\usepackage{array}
\usepackage{geometry}
\usepackage{fancyhdr}
\usepackage{tcolorbox}
\usepackage{tikz}
\usetikzlibrary{shapes,arrows,positioning,fit,backgrounds}

\geometry{margin=1in}

% Colors
\definecolor{qnkpurple}{RGB}{128,0,128}
\definecolor{qnkblue}{RGB}{0,100,180}
\definecolor{codebg}{RGB}{245,245,245}
\definecolor{codeframe}{RGB}{200,200,200}

% Listings style
\lstset{
    backgroundcolor=\color{codebg},
    basicstyle=\ttfamily\small,
    breaklines=true,
    frame=single,
    frameround=tttt,
    rulecolor=\color{codeframe},
    numbers=left,
    numberstyle=\tiny\color{gray},
    keywordstyle=\color{qnkblue}\bfseries,
    commentstyle=\color{green!50!black},
    stringstyle=\color{orange},
    showstringspaces=false,
    tabsize=2
}

% Custom environments
\newtheorem{theorem}{Theorem}[section]
\newtheorem{lemma}[theorem]{Lemma}
\newtheorem{definition}[theorem]{Definition}

% Header/Footer
\pagestyle{fancy}
\fancyhf{}
\rhead{Q-NarwhalKnight Protocol}
\lhead{Post-Quantum Recursive SNARKs}
\rfoot{Page \thepage}

\title{
    \vspace{-1cm}
    \textbf{Eliminating Weak Subjectivity via\\Post-Quantum Recursive SNARKs}\\[0.5cm]
    \large Technical Review: Q-NarwhalKnight Cryptographic Light Client Protocol\\[0.3cm]
    \normalsize Version 1.0.0
}

\author{
    Q-NarwhalKnight Protocol Team\\
    \texttt{protocol@quillon.xyz}
}

\date{December 2024}

\begin{document}

\maketitle

\begin{abstract}
This document presents a novel approach to eliminating weak subjectivity in BFT consensus systems using \textbf{post-quantum recursive SNARKs}. We leverage Q-NarwhalKnight's existing LatticeGuard (RLWE-based zk-SNARK) and ZK-STARK infrastructure to create an \textbf{Incrementally Verifiable Computation (IVC)} chain that allows new nodes to cryptographically verify the entire blockchain history in constant time ($\sim$10ms) without trusting any checkpoint provider.

This is the first design for \textbf{post-quantum recursive proofs for hybrid DAG-BFT consensus}.
\end{abstract}

\tableofcontents
\newpage

\section{Problem Statement}

\subsection{What is Weak Subjectivity?}

In BFT/PoS consensus systems, new nodes joining the network face a fundamental problem. An attacker can create a fake chain with the same genesis but different history, and a new node cannot distinguish the real chain from the fake one without external trust.

\begin{tcolorbox}[colback=red!5!white,colframe=red!75!black,title=The Weak Subjectivity Problem]
\textbf{Real Chain:} $[G] \rightarrow [1] \rightarrow [2] \rightarrow \ldots \rightarrow [1000000]$\\
\textbf{Fake Chain:} $[G] \rightarrow [1'] \rightarrow [2'] \rightarrow \ldots \rightarrow [1000000']$\\[0.3cm]
Both chains start from the same genesis $G$, but contain different history. A new node cannot distinguish them without external trust!
\end{tcolorbox}

\textbf{Why BFT systems have this problem:}
\begin{itemize}
    \item Validator sets change over time
    \item Old validators may have unbonded and sold keys
    \item Attacker can buy old keys and sign alternate history
    \item No ``proof of work'' anchoring history to physics
\end{itemize}

\subsection{Why This Matters}

\begin{table}[h]
\centering
\begin{tabular}{lccc}
\toprule
\textbf{System} & \textbf{Bootstrap Trust} & \textbf{Verification Time} & \textbf{Post-Quantum} \\
\midrule
Bitcoin & None (verify from genesis) & Hours-Days & No \\
Ethereum 2.0 & Checkpoint trust & Minutes & No \\
Current Q-NarwhalKnight & Checkpoint trust & Minutes & Yes \\
\textbf{Proposed Q-NarwhalKnight} & \textbf{None (cryptographic)} & \textbf{$\sim$10ms} & \textbf{Yes} \\
\bottomrule
\end{tabular}
\caption{Comparison of blockchain bootstrap mechanisms}
\end{table}

\subsection{Goal}

Create a system where:
\begin{enumerate}
    \item New nodes verify entire chain history in \textbf{constant time} ($\sim$10ms)
    \item \textbf{No trusted checkpoints} -- purely cryptographic verification
    \item \textbf{Post-quantum secure} -- resistant to quantum attacks
    \item \textbf{Decentralized proof generation} -- no single prover
\end{enumerate}

\section{Background: Current Q-NarwhalKnight Architecture}

\subsection{Consensus Layers}

Q-NarwhalKnight implements a hybrid consensus with four layers:

\begin{enumerate}
    \item \textbf{Layer 1: Lightweight Mining} -- CPU-friendly proof-of-computation with memory-hard operations for Sybil resistance
    \item \textbf{Layer 2: VDF Leader Election} -- Genus-2 Jacobian VDF with sequential computation (no parallel speedup)
    \item \textbf{Layer 3: DAG Structure} -- Parallel block production with multiple parents per block
    \item \textbf{Layer 4: BFT Finality} -- 2f+1 validator signatures using SQIsign/Dilithium5 post-quantum signatures
\end{enumerate}

\subsection{Existing ZK Infrastructure}

Q-NarwhalKnight has three ZK systems:

\subsubsection{LatticeGuard (Post-Quantum SNARK)}

Based on Ring-LWE assumption with security levels PQ128, PQ192, PQ256:

\begin{lstlisting}[language=Rust,caption=LatticeGuard proof structure]
pub struct LatticeGuardProof {
    commitments: Vec<LatticeCommitment>,
    evaluations: (Scalar, Scalar, Scalar),
    product_proofs: Vec<ApproximateProductProof>,
    transcript_state: [u8; 32],
}
\end{lstlisting}

Key parameters: Dimension 1024-4096, Modulus 32-64 bits, Proof size 10-50 KB.

\subsubsection{ZK-STARK (Hash-Based, Inherently PQ)}

Transparent setup with GPU acceleration, targeting 50K+ TPS.

\subsubsection{Traditional SNARKs}

Groth16, PLONK, Marlin, Sonic for non-PQ applications where performance is critical.

\subsection{Network Layer (libp2p)}

Gossipsub topics for P2P communication:
\begin{itemize}
    \item \texttt{/qnk/testnet/blocks} -- Block propagation
    \item \texttt{/qnk/testnet/peer-heights} -- Height announcements
    \item \texttt{/qnk/testnet/bft-votes} -- BFT signature collection
\end{itemize}

\section{Solution Overview: Recursive Proof Chain}

\subsection{Core Idea: Incrementally Verifiable Computation (IVC)}

Instead of verifying each block individually, we create a \textbf{single proof} that attests to the validity of all blocks from genesis to current height.

\begin{definition}[Recursive Epoch Proof]
Let $\pi_n$ be the proof for epoch $n$. Then:
$$\pi_n = \text{Prove}\left(\text{Verify}(\pi_{n-1}) = 1, \text{blocks}_n, \text{sigs}_n, \text{state}_n\right)$$
\end{definition}

The key property is that verification time is $O(1)$ regardless of chain length.

\subsection{What Each Epoch Proof Contains}

\begin{lstlisting}[language=Rust,caption=Epoch proof public inputs]
pub struct EpochPublicInputs {
    pub previous_state_root: [u8; 32],
    pub current_state_root: [u8; 32],
    pub epoch: u64,
    pub height_range: (u64, u64),
    pub validator_set_hash: [u8; 32],
    pub signature_count: u32,
}
\end{lstlisting}

\subsection{The Recursive Circuit}

The epoch transition circuit verifies:
\begin{enumerate}
    \item $C_1$: Previous proof verification $\text{Verify}(\pi_{n-1}, \text{prev\_root}) = 1$
    \item $C_2$: Validator set hash correctness
    \item $C_3$: BFT threshold $\geq 2f+1$ valid signatures
    \item $C_4$: All epoch blocks are valid
    \item $C_5$: State transition correctness
    \item $C_6$: Merkle root computation
\end{enumerate}

\section{Circuit Designs}

\subsection{LatticeGuard Verifier Circuit (Recursive Component)}

The most critical component: a circuit that verifies a LatticeGuard proof inside itself.

\begin{theorem}[Recursive Verification Complexity]
The LatticeGuardVerifierCircuit requires approximately 100,000 R1CS constraints for verification of a proof with dimension 1024.
\end{theorem}

The circuit performs:
\begin{enumerate}
    \item \textbf{Commitment verification} -- Verify RLWE ciphertexts are well-formed
    \item \textbf{Fiat-Shamir transcript reconstruction} -- Recompute challenges using Poseidon hash
    \item \textbf{Polynomial evaluation verification} -- Check evaluations at challenge point
    \item \textbf{Approximate product verification} -- Verify R1CS satisfaction with bounded error
\end{enumerate}

\subsection{BFT Signature Verification Circuit}

Proves that $\geq 2f+1$ validators signed the epoch blocks using Dilithium5 signatures.

\begin{lstlisting}[language=Rust,caption=BFT signature circuit]
pub struct BFTSignatureCircuit {
    n_validators: usize,
    f: usize,  // Byzantine threshold
    validator_keys: Vec<DilithiumPublicKey>,
    signatures: Vec<Option<DilithiumSignature>>,
    message: [u8; 32],
}
\end{lstlisting}

Dilithium verification requires approximately 100,000 constraints per signature.

\subsection{State Transition Circuit}

Verifies that epoch state transitions are valid:
\begin{itemize}
    \item Block hash computation correctness
    \item DAG parent existence verification
    \item VDF output correctness (lightweight check)
    \item Transaction validity
    \item State update computation
\end{itemize}

\subsection{Complete Epoch Circuit}

\begin{table}[h]
\centering
\begin{tabular}{lrl}
\toprule
\textbf{Component} & \textbf{Constraints} & \textbf{Notes} \\
\midrule
LatticeGuard Verifier & $\sim$100,000 & Recursive verification \\
BFT Signatures (5 sigs) & $\sim$500,000 & Minimum viable \\
State Transition & $\sim$200,000 & Depends on epoch size \\
Overhead & $\sim$50,000 & Glue logic \\
\midrule
\textbf{Total per Epoch} & $\sim$850,000 & Conservative estimate \\
\bottomrule
\end{tabular}
\caption{Epoch transition circuit constraint breakdown}
\end{table}

\section{Decentralized Proof Generation via libp2p}

\subsection{The Decentralization Challenge}

Generating recursive proofs is computationally expensive:
\begin{itemize}
    \item $\sim$850K constraints per epoch
    \item $\sim$30-60 seconds proving time (CPU)
    \item $\sim$5-10 seconds with GPU acceleration
\end{itemize}

We cannot rely on a single prover -- this would centralize trust.

\subsection{Decentralized Architecture}

Multiple prover nodes compete to generate epoch proofs:

\begin{enumerate}
    \item \textbf{Epoch Finalized} -- BFT consensus completes
    \item \textbf{Task Broadcast} -- Gossipsub: \texttt{/qnk/epoch-proof-task}
    \item \textbf{Parallel Proving} -- Multiple provers race
    \item \textbf{First Valid Wins} -- Gossipsub: \texttt{/qnk/epoch-proofs}
    \item \textbf{All Verify} -- 10ms verification by all nodes
    \item \textbf{DHT Storage} -- Key: \texttt{/qnk/proofs/epoch/\{N\}}
\end{enumerate}

\subsection{libp2p Protocol Specification}

New gossipsub topics:
\begin{lstlisting}[language=Rust]
pub const TOPIC_EPOCH_PROOF_TASK: &str = "/qnk/epoch-proof-task";
pub const TOPIC_EPOCH_PROOFS: &str = "/qnk/epoch-proofs";
pub const TOPIC_PROOF_VERIFICATION: &str = "/qnk/proof-verification";
\end{lstlisting}

\subsection{Incentive Mechanism}

\begin{lstlisting}[language=Rust,caption=Proof reward calculation]
pub fn calculate_reward(submission: &EpochProofSubmission, task: &EpochProofTask) -> u64 {
    let mut reward = BASE_REWARD;

    // Speed bonus for fast proofs
    if submission.proving_time_ms < TARGET_PROVING_TIME_MS {
        let speedup = TARGET_PROVING_TIME_MS - submission.proving_time_ms;
        reward += SPEED_BONUS * speedup / TARGET_PROVING_TIME_MS;
    }

    // Late penalty (no negative rewards)
    if now > task.deadline {
        let penalty = LATE_PENALTY_PER_SECOND * (now - task.deadline);
        reward = reward.saturating_sub(penalty);
    }

    reward
}
\end{lstlisting}

\subsection{Light Client Sync Protocol}

The key function -- trustless bootstrap in $\sim$10ms:

\begin{lstlisting}[language=Rust,caption=Light client bootstrap]
pub async fn bootstrap(&mut self) -> Result<()> {
    // Request proof from multiple peers
    let responses = self.request_from_multiple_peers(request).await?;
    let best_response = self.find_consensus_response(&responses)?;

    // CRITICAL: Verify the proof - trustless!
    let is_valid = self.verifier.verify(
        &best_response.proof,
        &public_inputs,
    )?;

    // Verification time: ~10ms
    // Trust required: NONE
}
\end{lstlisting}

\section{Implementation Roadmap}

\begin{enumerate}
    \item \textbf{Phase 1: Circuit Foundations (4-6 weeks)}
    \begin{itemize}
        \item Poseidon hash gadget for LatticeGuard
        \item Dilithium signature verification circuit
        \item Merkle tree verification circuit
    \end{itemize}

    \item \textbf{Phase 2: Recursive Prover (6-8 weeks)}
    \begin{itemize}
        \item LatticeGuardVerifierCircuit
        \item BFTSignatureCircuit
        \item StateTransitionCircuit
        \item EpochTransitionCircuit
    \end{itemize}

    \item \textbf{Phase 3: P2P Integration (4-6 weeks)}
    \begin{itemize}
        \item New gossipsub topics
        \item ProverNode implementation
        \item Proof verification and storage
    \end{itemize}

    \item \textbf{Phase 4: Light Client (3-4 weeks)}
    \begin{itemize}
        \item LightClient bootstrap
        \item Proof request/response protocol
        \item Wallet integration
    \end{itemize}

    \item \textbf{Phase 5: Optimization (Ongoing)}
    \begin{itemize}
        \item GPU acceleration
        \item Signature aggregation
        \item Proof compression
    \end{itemize}
\end{enumerate}

\section{Security Analysis}

\subsection{Threat Model}

\begin{table}[h]
\centering
\begin{tabular}{p{4cm}p{8cm}}
\toprule
\textbf{Threat} & \textbf{Mitigation} \\
\midrule
Malicious prover generates invalid proof & All nodes verify proofs before accepting \\
Prover collusion & Multiple independent provers race; any valid proof accepted \\
Proof withholding & Timeout triggers re-proving by other nodes \\
Long-range attack (fake history) & Recursive proof verifies entire history cryptographically \\
Quantum attack on RLWE & Parameters chosen for 128+ bit post-quantum security \\
\bottomrule
\end{tabular}
\caption{Threat model and mitigations}
\end{table}

\subsection{Security Assumptions}

\begin{enumerate}
    \item \textbf{RLWE hardness}: Ring-LWE problem is hard for quantum computers
    \item \textbf{Hash collision resistance}: BLAKE3/Poseidon are collision-resistant
    \item \textbf{Honest majority for BFT}: $<33\%$ Byzantine validators
    \item \textbf{At least one honest prover}: Some prover generates valid proofs
\end{enumerate}

\section{Performance Projections}

\subsection{Proving Times}

\begin{table}[h]
\centering
\begin{tabular}{lcc}
\toprule
\textbf{Hardware} & \textbf{Estimated Time} & \textbf{Speedup} \\
\midrule
CPU (16 cores) & 45-60 seconds & 1x (baseline) \\
GPU (RTX 4090) & 8-12 seconds & 5-6x \\
GPU Cluster & 3-5 seconds & 10-15x \\
FPGA (future) & $<$1 second & 50x+ \\
\bottomrule
\end{tabular}
\caption{Proving time estimates by hardware}
\end{table}

\subsection{Proof Sizes and Verification}

\begin{itemize}
    \item \textbf{Total Light Client Proof}: $\sim$50 KB
    \item \textbf{Verification Time}: $\sim$10-20 ms
    \item \textbf{Verification Complexity}: $O(1)$ regardless of chain length
\end{itemize}

\section{Comparison with Existing Work}

\subsection{Mina Protocol}

\begin{table}[h]
\centering
\begin{tabular}{lcc}
\toprule
\textbf{Aspect} & \textbf{Mina} & \textbf{Q-NarwhalKnight} \\
\midrule
Proof System & Pickles (Pasta curves) & LatticeGuard (RLWE) \\
Post-Quantum & No & Yes \\
Consensus & Ouroboros Samasika & DAG-BFT + VDF Mining \\
Proof Size & $\sim$22 KB & $\sim$50 KB \\
Verification & $\sim$1 second & $\sim$10 ms \\
\bottomrule
\end{tabular}
\caption{Comparison with Mina Protocol}
\end{table}

\subsection{Key Innovation}

\textbf{Q-NarwhalKnight is the first to combine:}
\begin{enumerate}
    \item Post-quantum recursive proofs (LatticeGuard)
    \item Hybrid consensus (VDF mining + BFT)
    \item Decentralized proving via libp2p
    \item Elimination of weak subjectivity for BFT
\end{enumerate}

\section{Open Research Questions}

\subsection{Efficiency Improvements}
\begin{itemize}
    \item Can we aggregate Dilithium signatures to reduce BFT circuit size?
    \item Can we prove sub-epochs and combine (incremental proving)?
    \item Can recursive proof size be reduced below 50KB?
\end{itemize}

\subsection{Security Questions}
\begin{itemize}
    \item What RLWE parameters balance security vs. performance?
    \item What if all provers go offline?
    \item How long should proofs be valid?
\end{itemize}

\subsection{Economic Questions}
\begin{itemize}
    \item What reward structure ensures sufficient provers?
    \item Should there be a proof marketplace?
    \item Should validators be required to prove?
\end{itemize}

\section{Conclusion}

This design eliminates weak subjectivity from Q-NarwhalKnight's BFT consensus layer using post-quantum recursive SNARKs. The key innovations are:

\begin{enumerate}
    \item \textbf{LatticeGuard Recursion}: First recursive proof system based on RLWE
    \item \textbf{Decentralized Proving}: P2P network of competing provers via libp2p
    \item \textbf{Constant-Time Verification}: New nodes verify entire history in $\sim$10ms
    \item \textbf{Full Post-Quantum Security}: All cryptographic components are quantum-resistant
\end{enumerate}

This represents a significant step forward in blockchain technology: \textbf{trustless light clients for BFT consensus} without any social-layer assumptions.

\appendix

\section{Glossary}

\begin{table}[h]
\centering
\begin{tabular}{lp{10cm}}
\toprule
\textbf{Term} & \textbf{Definition} \\
\midrule
IVC & Incrementally Verifiable Computation -- proofs that verify other proofs \\
RLWE & Ring Learning With Errors -- post-quantum hardness assumption \\
R1CS & Rank-1 Constraint System -- arithmetic circuit representation \\
Weak Subjectivity & Need for trusted checkpoints in BFT systems \\
BFT & Byzantine Fault Tolerance -- consensus despite malicious actors \\
VDF & Verifiable Delay Function -- sequential computation proof \\
\bottomrule
\end{tabular}
\end{table}

\end{document}
